% Esquisse d'un programme (cote 119, copy 1): archive cote 119, site page 39 (PDF page 40), typescript page 37.
% typescript. Draft compared: EsquisseFr.pdf (text layer):1340-1404.
% AI-transcribed by Claude Opus 5.5 (claude-opus-5-5, adaptive thinking, effort high), Message Batches API, batch esq1 (msgbatch_01LNsogFyF3y6KSp9UbtJtBk), 2026-11-12; not yet checked by eye.
\scanpage{119/39}
\mspage{37}

6. Un des théorèmes de fondements de topologie (modérée)
les plus intéressants qu'il faudrait développer, serait un théo-%
rème de ``dévissage'' (encore !) d'une application modérée propre
d'espace modérés,

\[ f : X \longrightarrow Y , \]

via une filtration décroissante de $Y$ par des sous-espaces mo-%
dérés fermés $Y^{i}$, tels que au-dessus des ``strates ouvertes''
$Y^{i} \setminus Y^{i-1}$ de cette filtration, $f$ induise une fibration locale-%
ment triviale (du point de vue modéré, il va sans dire). Un tel
énoncé devrait encore se généraliser et se préciser de diverses
façons, notamment en demandant l'existence d'un dévissage ana-%
logue \uline{simultané}, pour $X$ et une famille finie donnée de sous-%
espaces (modérés) fermés de $X$. Egalement la notion même de fi-%
bration localement triviale au sens modéré peut se renforcer
considérablement, en tenant compte du fait que les strates ou-%
vertes $U_{i}$ sont \uline{mieux} que des espaces à structure modérée pu-%
rement locale, du fait qu'ils sont obtenus comme différence de
deux espaces modérés, compacts si $Y$ était compact. Entre la
notion d'espace modéré compact (qui se réalise comme un des
``modèles'' de départ dans un $\mathbb{R}^{n}$) et celle d'espace ``localement
modéré'' (localement compact) qui s'en déduit de façon assez évi-%
dente, il y a une notion un peu plus délicate d'\uline{espace} \uline{``globa-}%
\uline{lement modéré}'' $X$, obtenu comme différence $X^{\wedge}\setminus Y$ de deux espaces
modérés compacts, étant entendu qu'on ne distingue pas entre
l'espace défini par une paire $(X^{\wedge},Y)$, et celui défini par une
paire $(X'^{\wedge},Y')$ qui s'en déduit par une application modérée
(nécessairement propre)

\[ g : X'^{\wedge} \longrightarrow X^{\wedge} \]

induisant une bijection $g^{-1}(X) \xrightarrow{\ \sim\ } X$, en prenant $Y' = g^{-1}(Y)$.
L'exemple naturel le plus intéressant peut-être est celui où on
part d'un schéma séparé de type fini sur $\mathbb{C}$ ou sur $\mathbb{R}$, en prenant

\margin{39}
